\chapter{Materie uit voedsel}\label{ch:g6:matter-from-food}

Een kalf van \qty{40}{kg} wordt in twee jaar een koe van \qty{600}{kg}.
Een halve ton nieuwe koe --- vlees, \omterm{def:g3:bones-and-muscles:skeleton}{bot}, \omterm{def:g1:the-five-senses:sense}{huid}, hoorn --- is ergens
vandaan gekomen. De koe at gras; de nieuwe materie \emph{is geen gras};
en toch zou zonder het gras niets ervan bestaan. Wat \omterm{def:g1:animals-around-us:animal}{dieren} met hun
voedsel doen is het onderwerp van dit hoofdstuk --- de kant van de
\omterm{def:g5:food-webs:roles}{consumenten} in de boekhouding van \cref{ch:g6:plants-are-producers}.

\section{Dieren bouwen hun materie uit voedsel}

\begin{proposition}[Consumenten zetten organische materie om]\label{prop:g6:matter-from-food:transform}
Een \omterm{def:g1:animals-around-us:animal}{dier} bouwt zijn eigen \omterm{def:g6:plants-are-producers:matter}{organische materie} --- \omterm{def:g3:bones-and-muscles:muscle}{spier}, \omterm{def:g3:bones-and-muscles:skeleton}{bot}, vet, \omterm{def:g1:animals-around-us:covering}{veer},
\omterm{def:g1:animals-around-us:covering}{schelp} --- uit de \emph{\omterm{def:g6:plants-are-producers:matter}{organische materie} die het eet}. De \omterm{def:g4:journey-of-food:digestion}{vertering}
(\cref{def:g4:journey-of-food:digestion}) haalt het voedsel uit elkaar
tot \omterm{def:g4:journey-of-food:digestion}{voedingsstoffen}; het \omterm{prop:g4:journey-of-food:absorb}{bloed} bezorgt ze
(\cref{prop:g4:heart-and-blood:carries}); en de cellen van het lichaam
(\cref{prop:g5:first-look-at-cells:growth}) zetten ze weer in elkaar tot
de eigen stof van het \omterm{def:g1:animals-around-us:animal}{dier}. \omterm{def:g1:animals-around-us:animal}{Dieren} zijn \emph{omzetters} van materie,
nooit makers-uit-mineralen: dat vak hoort alleen de \omterm{def:g5:food-webs:roles}{producenten} toe.
\end{proposition}
\admitted

\begin{example}[Gras wordt koe]\label{ex:g6:matter-from-food:cow}
Volg die halve ton: de koe graast uur na uur
(\cref{prop:g2:what-animals-eat:daily} verklaarde die lange uren); haar
zware \omterm{def:g4:journey-of-food:digestion}{vertering} (\cref{exo:g4:journey-of-food:11}) haalt het gras uit
elkaar tot \omterm{def:g4:journey-of-food:digestion}{voedingsstoffen}; haar \omterm{prop:g4:journey-of-food:absorb}{bloed} brengt die naar uier, \omterm{def:g3:bones-and-muscles:muscle}{spier} en het
groeiende kalf in haar; haar cellen bouwen koeienmaterie uit
grasmaterie. De weide loopt, heel letterlijk, in een nieuwe vorm rond.
\end{example}

\begin{remark}[Materie wisselt van eigenaar, niet van aard]\label{rem:g6:matter-from-food:hands}
Bij elke pijl van een \omterm{def:g3:food-chains:chain}{voedselketen}
(\cref{def:g3:food-chains:chain}) wisselt \omterm{def:g6:plants-are-producers:matter}{organische materie} van eigenaar
en wordt ze herbouwd --- grasmaterie tot koeienmaterie, koeienmaterie (de
\omterm{def:g2:animals-grow-up:born}{melk}, het vlees) tot kindermaterie. Het \emph{organische} komt nooit uit
het niets: \omterm{def:g5:food-webs:roles}{consumenten} geven het door en vervormen het. Alleen aan het
groene begin van de keten werd het nieuw gemaakt
(\cref{prop:g6:plants-are-producers:firstlink}); alleen aan het einde bij
de \omterm{def:g5:food-webs:roles}{afbrekers}, in het volgende hoofdstuk, keert het terug naar mineraal.
\end{remark}

\section{Groei is materieproductie}

\begin{definition}[De materieproductie van een dier]\label{def:g6:matter-from-food:production}
De \emph{materieproductie}\index{materieproductie} van een \omterm{def:g1:animals-around-us:animal}{dier} is de
nieuwe \omterm{def:g6:plants-are-producers:matter}{organische materie} die het bouwt: de groei van het jong (de
kilo's die het kalf erbij krijgt), jaarlijkse vernieuwingen (\omterm{def:g1:animals-around-us:covering}{vacht} die
verharen en weer aangroeien, geweien, \omterm{def:g1:animals-around-us:covering}{veren}), herstel (het geheelde \omterm{def:g3:bones-and-muscles:skeleton}{bot}
uit \cref{rem:g3:bones-and-muscles:alive}) en producten (\omterm{def:g2:animals-grow-up:born}{melk}, eieren,
wol). Het is allemaal herbouwd voedsel.
\end{definition}

\begin{example}[Productie wegen]\label{ex:g6:matter-from-food:weighing}
Boeren wegen de productie dagelijks: een melkkoe geeft zo'n
\qty{25}{kg} \omterm{def:g2:animals-grow-up:born}{melk} per dag --- \omterm{def:g6:plants-are-producers:matter}{organische materie} die ze dag na dag uit
voer moet herbouwen; het \omterm{def:g2:animals-grow-up:born}{ei} van een kip is zo'n \qty{60}{g} ingepakte
proviand (\cref{rem:g2:animals-grow-up:egg}), 's nachts uit haar graan
in elkaar gezet. Geen voer, geen productie: de boekhouding klopt
genadeloos.
\end{example}

\begin{proposition}[Niet al het voedsel wordt lichaam]\label{prop:g6:matter-from-food:losses}
Een \omterm{def:g5:food-webs:roles}{consument} zet nooit al zijn voedsel om in nieuwe materie:
\begin{itemize}
  \item een deel wordt nooit opgenomen en verlaat het lichaam als
        ontlasting (\cref{prop:g4:journey-of-food:absorb});
  \item een groot deel wordt \emph{verbrand} --- verbonden met de
        zuurstof uit de adem
        (\cref{prop:g4:breathing:exchange}) --- om beweging en warmte
        aan te drijven: materie die als brandstof uitgegeven wordt en
        als uitgeademd gas vertrekt;
  \item alleen de rest wordt groei en producten.
\end{itemize}
Van het gras dat een koe eet eindigt maar een bescheiden deel als nieuwe
koe --- en daarom is elke verdieping van de piramide uit
\cref{prop:g5:food-webs:pyramid} zoveel kleiner dan die eronder.
\end{proposition}
\admitted

\begin{example}[De piramide, eindelijk verklaard]\label{ex:g6:matter-from-food:pyramid}
De tonnen gras van een wei voeden een kudde konijnen; de kilo's van de
konijnen voeden een vos of twee. Bij elke pijl wordt de meeste materie
voor het leven verbrand of gaat ze verloren, en wordt maar een fractie
als lichaam naar boven doorgegeven. Grote jagers zijn zeldzaam
(\cref{prop:g5:food-webs:pyramid}) omdat ze van de \emph{restanten van
restanten} leven --- rekenkunde, geen pech.
\end{example}

\section{Menu's, tanden en hetzelfde ene vak}

\begin{proposition}[Verschillende menu's, dezelfde omzetting]\label{prop:g6:matter-from-food:menus}
\omterm{def:g2:what-animals-eat:kinds}{Planteneter}, \omterm{def:g2:what-animals-eat:kinds}{vleeseter} en \omterm{def:g2:what-animals-eat:kinds}{alleseter}
(\cref{def:g2:what-animals-eat:kinds}) benoemen verschillende
\emph{bronnen} van \omterm{def:g6:plants-are-producers:matter}{organische materie}, geen verschillende vakken. De koe
die gras herbouwt en de vos die konijn herbouwt draaien allebei de
omzetting van \cref{prop:g6:matter-from-food:transform} --- met een
uitrusting die bij het menu past
(\cref{prop:g2:what-animals-eat:match},
\cref{exo:g4:journey-of-food:11}): gras vraagt maalkiezen en een enorme
\omterm{def:g4:journey-of-food:digestion}{vertering}, vlees beloont grijpers en een korte darm.
\end{proposition}
\admitted

\begin{method}[De materie van een dier controleren]\label{met:g6:matter-from-food:audit}
Stel voor elk \omterm{def:g1:animals-around-us:animal}{dier} zijn materiebalans op:
\begin{enumerate}
  \item \emph{in}: welke \omterm{def:g6:plants-are-producers:matter}{organische materie} eet het, en van wie? (zijn
        binnenkomende pijlen in het \omterm{def:g5:food-webs:web}{voedselweb},
        \cref{def:g5:food-webs:web});
  \item \emph{gebouwd}: welke groei, vernieuwingen en producten maakt
        het?
  \item \emph{uit}: ontlasting, en materie die voor beweging en warmte
        verbrand is;
  \item controleer de balans: de productie moet de bescheiden rest zijn
        --- als een bewering haar groter maakt (``deze vijver levert
        meer vismaterie op dan de vijver aan voedsel groeit''), dan
        klopt de bewering niet.
\end{enumerate}
\end{method}

\begin{example}[De controle in de praktijk]\label{ex:g6:matter-from-food:audituse}
Controleer de merel (\cref{ex:g2:what-animals-eat:winter}): in ---
wormen, \omterm{prop:g3:sorting-living-things:groups}{insecten}, bessen; gebouwd --- zijn eigen groei, elk jaar nieuwe
\omterm{def:g1:animals-around-us:covering}{veren}, het legsel eieren; uit --- poep op de muur, en het grote verbrande
deel dat een jaar vliegen en zingen aandreef. Het dagelijkse
voedselgewicht van een zangvogel is enorm voor zijn formaat, juist omdat
het vliegen zo'n groot deel van de boekhouding opbrandt.
\end{example}

\section{Oefeningen}

\begin{exercise}[$\star$]\label{exo:g6:matter-from-food:1}
Beschrijf het vak van de \omterm{def:g5:food-webs:roles}{consumenten}, en zet het in één zin per stuk
tegenover dat van de \omterm{def:g5:food-webs:roles}{producenten}.
\end{exercise}

\begin{exercise}[$\star$]\label{exo:g6:matter-from-food:2}
Noem de stappen die gras tot koeienspier maken, met bij elke stap het
hoofdstuk dat haar behandelde.
\end{exercise}

\begin{exercise}[$\star$]\label{exo:g6:matter-from-food:3}
Geef vier vormen van \omterm{def:g6:matter-from-food:production}{materieproductie} bij een \omterm{def:g1:animals-around-us:animal}{dier}, met van elk een
voorbeeld.
\end{exercise}

\begin{exercise}[$\star$]\label{exo:g6:matter-from-food:4}
Noem de drie bestemmingen van het voedsel van een \omterm{def:g5:food-webs:roles}{consument}, volgens
\cref{prop:g6:matter-from-food:losses}.
\end{exercise}

\begin{exercise}[$\star$]\label{exo:g6:matter-from-food:5}
Waarom heeft een melkkoe elke dag opnieuw voer nodig om \omterm{def:g2:animals-grow-up:born}{melk} te blijven
geven?
\end{exercise}

\begin{exercise}[$\star$]\label{exo:g6:matter-from-food:6}
Een \omterm{def:g2:animals-grow-up:born}{ei} bevat \qty{60}{g} \omterm{def:g6:plants-are-producers:matter}{organische materie}. Waar is die gefabriceerd ---
in de kip, of daarvóór? Ontwar het zorgvuldig.
\end{exercise}

\begin{exercise}[$\star\star$]\label{exo:g6:matter-from-food:7}
Leg met \cref{prop:g6:matter-from-food:losses} uit waarom de
voedselpiramide bij elke verdieping smaller wordt.
\end{exercise}

\begin{exercise}[$\star\star$]\label{exo:g6:matter-from-food:8}
Laat \cref{met:g6:matter-from-food:audit} lopen over een huiskat: in,
gebouwd, uit --- met minstens twee posten per regel.
\end{exercise}

\begin{exercise}[$\star\star$]\label{exo:g6:matter-from-food:9}
``De vos en de koe doen verschillende dingen met voedsel.'' Verbeter die
zin met \cref{prop:g6:matter-from-food:menus} --- wat verschilt er, en
wat is identiek?
\end{exercise}

\begin{exercise}[$\star\star$]\label{exo:g6:matter-from-food:10}
Een kind groeit in een jaar \qty{3}{kg} terwijl het enkele honderden
kilo's voedsel eet. Waar bleef de rest? Antwoord met de boekhouding.
\end{exercise}

\begin{exercise}[$\star\star$]\label{exo:g6:matter-from-food:11}
Het huisje van de slak is minerale stof, en toch telde
\cref{ex:g6:plants-are-producers:sorting} het als maaksel van de slak.
Breng dat met elkaar in overeenstemming: wat deed het lichaam van de
slak, en uit welke grondstoffen, om het te bouwen?
\end{exercise}

\begin{exercise}[$\star\star\star$]\label{exo:g6:matter-from-food:12}
Eén hectare land kan meer mensen voeden met aardappelen dan met varkens
die op die aardappelen gehouden worden. Leg dat uit met
\cref{prop:g6:matter-from-food:losses} --- waar gaat het verschil heen
als de aardappelen de omweg via het varken nemen?
\end{exercise}

\section{Probleem: De boekhouding van de boerderij}

\begin{problem}[{Weekendprobleem --- een kleine boerderij doorgelicht,
van akker tot tafel}]\label{pb:g6:matter-from-food:1}
Het boerderijtje van oma: een weiland met één koe, een kippenhok met tien
kippen die graan krijgen, een moestuin en de keuken van het gezin. Dit
weekend houd jij de materieboekhouding bij.

\medskip\noindent\textbf{Deel I --- De verdieping van de \omterm{def:g5:food-webs:roles}{producenten}.}
\begin{enumerate}
  \item Noem de \omterm{def:g5:food-webs:roles}{producenten} van het bedrijf. Wat zijn hun minerale
        grondstoffen, en wat is hun energiebron?
  \item Het weiland laat in een seizoen ongeveer \qty{10000}{kg}
        grasmaterie groeien. In wier \omterm{def:g1:plants-around-us:parts}{bladeren} is elke gram daarvan
        gefabriceerd?
  \item Oma strooit de mest uit het kippenhok over de moestuin. Welke
        regel uit \cref{prop:g5:people-and-nature:lasting} houdt ze aan,
        en wat vult de mest precies aan?
  \item Waarom kun je de koe niet gewoon rechtstreeks \omterm{prop:g4:what-plants-need:needs}{minerale zouten} en
        water voeren en het weiland uitsparen?
\end{enumerate}

\medskip\noindent\textbf{Deel II --- De verdieping van de \omterm{def:g5:food-webs:roles}{consumenten}.}
\begin{enumerate}[resume]
  \item De koe eet ruwweg \qty{70}{kg} gras per dag en geeft ongeveer
        \qty{20}{kg} \omterm{def:g2:animals-grow-up:born}{melk}. Noem de drie bestemmingen van de ontbrekende
        \qty{50}{kg}.
  \item De kippen maken van graan eieren. Schrijf de boekhouding van een
        kip op --- in, gebouwd, uit --- en zeg waarom een ruiende kip
        (die al haar \omterm{def:g1:animals-around-us:covering}{veren} vernieuwt) een tijdlang minder eieren legt.
  \item Een buizerd pakt een kuiken. Breid de \omterm{def:g3:food-chains:chain}{voedselketen} van het
        bedrijf tot hem uit, en zeg bij welke schakels materie
        \emph{gemaakt} werd en bij welke ze alleen omgezet werd.
  \item Het gezin eet eieren, \omterm{def:g2:animals-grow-up:born}{melk} en groente. Volg de \omterm{def:g6:plants-are-producers:matter}{organische materie}
        van één ontbijt --- brood, gekookt \omterm{def:g2:animals-grow-up:born}{ei}, \omterm{def:g2:animals-grow-up:born}{melk} --- elk terug naar
        het \omterm{def:g1:plants-around-us:parts}{blad} waar ze het eerst gefabriceerd is.
\end{enumerate}

\medskip\noindent\textbf{Deel III --- De rekenkunde van de
verdiepingen.}
\begin{enumerate}[resume]
  \item Oma zegt: ``het weiland voedt één koe; het zou hooguit een klein
        vossenhol voeden als vossen er van gras gevoede konijnen aten''.
        Onderbouw haar schatting met de verliezen op elke verdieping.
  \item De moestuin voedt het gezin rechtstreeks. Leg met de redenering
        van \cref{exo:g6:matter-from-food:12} uit waarom de moestuin per
        vierkante meter meer mensen voedt dan de route van weiland naar
        koe naar \omterm{def:g2:animals-grow-up:born}{melk}.
  \item Bij droogte stopt het gras met groeien. Voorspel op volgorde wat
        er met de \omterm{def:g2:animals-grow-up:born}{melk} van de koe, met haar gewicht en met de regel
        ``gebouwd'' in de boekhouding gebeurt --- en waarom geen extra
        zonneschijn kan helpen zolang de regen uitblijft
        (\cref{pb:g6:plants-are-producers:1}, vraag 7, heeft de sleutel).
  \item Sluit de boekhouding af met één zin per verdieping: \omterm{def:g5:food-webs:roles}{producenten},
        \omterm{def:g5:food-webs:roles}{consumenten} --- en noem wie er in de boekhouding nog ontbreekt
        (het onderwerp van het volgende hoofdstuk).
\end{enumerate}
\end{problem}
